% ============================================================
%  A Process-Based Coherence Field Bridging Classical and
%  Quantum Dynamics (Not-Math Framework)
%  Version 2.0 — Verified and Corrected
%  Overleaf-compatible LaTeX
% ============================================================
\documentclass[11pt, a4paper]{article}

\usepackage[margin=2.2cm]{geometry}
\usepackage{amsmath, amssymb, amsthm, mathtools}
\usepackage{booktabs, array, longtable}
\usepackage{xcolor}
\usepackage{hyperref}
\usepackage{titlesec}
\usepackage{enumitem}
\usepackage{fancyhdr}
\usepackage{setspace}

\hypersetup{
  colorlinks=true,
  linkcolor=blue!50!black,
  urlcolor=blue!60!black,
}

\definecolor{nm-dark}{RGB}{30, 30, 30}
\definecolor{nm-gray}{RGB}{100, 100, 100}

\pagestyle{fancy}
\fancyhf{}
\rhead{\textcolor{nm-gray}{\small Not-Math Framework v2.0}}
\lhead{\textcolor{nm-gray}{\small Confidential --- Internal Use Only}}
\cfoot{\thepage}

\titleformat{\section}{\large\bfseries\color{nm-dark}}{\thesection}{0.6em}{}
\titleformat{\subsection}{\normalsize\bfseries\color{nm-dark}}{\thesubsection}{0.6em}{}

\setlength{\parindent}{0pt}
\setlength{\parskip}{6pt}
\onehalfspacing

\newtheorem{definition}{Definition}[section]
\newtheorem{axiom}{Axiom}[section]
\newtheorem{theorem}{Theorem}[section]
\newtheorem{proposition}{Proposition}[section]
\newtheorem{corollary}{Corollary}[section]
\newtheorem{remark}{Remark}[section]

% ============================================================
\begin{document}

\begin{center}
  {\LARGE \textbf{A Process-Based Coherence Field Bridging}}\\[4pt]
  {\LARGE \textbf{Classical and Quantum Dynamics}}\\[6pt]
  {\large (Not-Math Framework)}\\[8pt]
  {\color{nm-gray}\small Version 2.0 --- Verified and Corrected}\\[4pt]
  {\color{nm-gray}\small Joseph Forrester \quad|\quad Independent Researcher}\\[4pt]
  {\color{nm-gray}\small April 2026 \quad|\quad Confidential --- Internal Use Only}
\end{center}

\vspace{1em}
\hrule
\vspace{1em}

% ============================================================
\section*{Abstract}

We present a conservative field framework --- the Universal Field of Coherent
Processes (UFCP) --- in which conventional ``numbers'' are reinterpreted as
processual quantities that carry local tempo, phase, and coherent potential.
From a single Hermitian action with symmetric couplings we derive equations
of motion, Noether currents, and precise correspondence limits that reproduce
(i)~the Schr\"odinger equation and Born's rule, (ii)~Newtonian mechanics via
Ehrenfest/collective-coordinate limits, and (iii)~relativistic wave dynamics
with microcausality in Klein--Gordon/Dirac corners.

Measurement probabilities emerge by extremizing captured coherent potential
under a conserved norm, uniquely yielding the overlap-squared rule under
unitary invariance and orthogonal additivity. The framework makes falsifiable
predictions including a coupling-induced tempo dilation of coherently
``crowded'' oscillators and a coherence-resource trade-off in multi-path/party
interferometry.

We outline laboratory protocols and list explicit failure conditions that
would falsify the approach. This unifies classical, quantum, and relativistic
behaviors as projections of a single process-based dynamics while remaining
testable in precision metrology.

\vspace{0.5em}
\noindent\textbf{v2.0 Corrections:} Measurement derivation clarified as
axiom trade (not axiom elimination). Interaction kernel symmetry grounded
in CPT invariance. Classical limit decoherence mechanism specified. CITD
magnitude estimated for current instruments. Gravitational sector
extended beyond hand-wave. All derivations independently verified.

\tableofcontents

% ============================================================
\section{Introduction}

We develop a process-based field description where quantities are not static
scalars but coherent processes carrying a local tempo $\omega$, phase $\phi$,
and coherent-potential density $\rho$.

The ambition is pragmatic: recover known laws (classical, quantum,
relativistic) from one conservative action while exposing new, falsifiable
effects in precision metrology. Our construction is conservative in the
Noether sense (global $U(1)$, Poincar\'e, and gauge symmetries imply
conserved currents) and reduces to familiar theories in appropriate limits.

\subsection{Design Principles}

\begin{enumerate}[label=\arabic*.]
  \item \textbf{Single action.} All dynamics derive from one action
        functional $S$. No separate postulates for different regimes.
  \item \textbf{Hermitian operators.} The kinetic operator $\hat{\Omega}_g$
        is Hermitian, ensuring real energy eigenvalues and unitary time
        evolution.
  \item \textbf{Symmetric interactions.} The interaction kernel $W$ is
        symmetric, ensuring norm conservation and reciprocal coupling.
  \item \textbf{Falsifiability.} The framework makes specific, quantitative
        predictions that can be tested with current laboratory equipment.
  \item \textbf{Correspondence.} Known physics (Schr\"odinger, Newton,
        Klein--Gordon, Dirac) must emerge as exact limits, not approximations.
\end{enumerate}

% ============================================================
\section{Coherence Field and Action}

\begin{definition}[Coherence Field]
Let $\psi = \sqrt{\rho}\, e^{i\phi}$ be the coherence field on a
Lorentzian manifold $(M, g_{\mu\nu})$ with minimal gauge coupling
$D_\mu = \partial_\mu + i\frac{q}{\hbar}A_\mu$.

The field carries three physical attributes:
\begin{itemize}[nosep]
  \item $\rho(\mathbf{x}, t) \geq 0$: coherent-potential density (structural weight)
  \item $\phi(\mathbf{x}, t) \in \mathbb{R}$: phase (structural relationship)
  \item $\omega = \partial_t \phi$: local tempo (structural frequency)
\end{itemize}
\end{definition}

\begin{axiom}[UFCP Action]
The dynamics of the coherence field are governed by the action:
\begin{equation}
  \boxed{
  S = \int d^4x\, \sqrt{-g} \left[
    \frac{i\hbar}{2}\!\left(\psi^* D_t\psi - \psi D_t\psi^*\right)
    - \psi^* \hat{\Omega}_g \psi
    - U(\rho)
    - \frac{1}{2}\int \rho(\mathbf{x})\, W(\mathbf{x}-\mathbf{x}')\,
      \rho(\mathbf{x}')\, \sqrt{-g'}\, d^4x'
  \right]
  }
  \label{eq:action}
\end{equation}
where:
\begin{itemize}[nosep]
  \item $\hat{\Omega}_g$ is a Hermitian kinetic operator
        (e.g., $-\frac{\hbar^2}{2m}g^{ij}D_i D_j + V$ nonrelativistically)
  \item $U(\rho)$ is the local self-interaction potential (bounded below)
  \item $W(\mathbf{x}-\mathbf{x}')$ is the symmetric interaction kernel
\end{itemize}
\end{axiom}

\subsection{Equation of Motion}

Variation of $S$ with respect to $\psi^*$ yields:
\begin{equation}
  \boxed{
  i\hbar D_t \psi = \hat{\Omega}_g \psi
    + \frac{\partial U}{\partial \rho}\,\psi
    + (W * \rho)\,\psi
  }
  \label{eq:eom}
\end{equation}
where $(W * \rho)(\mathbf{x}) = \int W(\mathbf{x}-\mathbf{x}')\,
\rho(\mathbf{x}')\, \sqrt{-g'}\, d^4x'$.

\textbf{Verification:} The variation is standard for a Schr\"odinger-type
action. The $U(\rho)$ term contributes $\frac{\partial U}{\partial \rho}\psi$
because $\rho = |\psi|^2$ and $\frac{\delta \rho}{\delta \psi^*} = \psi$.
The $W$ term contributes $(W * \rho)\psi$ by the same chain rule. Both
terms are verified by direct computation.

\subsection{Conserved Currents (Noether)}

Global $U(1)$ symmetry ($\psi \to e^{i\alpha}\psi$) gives:
\begin{equation}
  \nabla_\mu J^\mu = 0, \qquad
  J^0 = \rho, \qquad
  \mathbf{J} = \rho\, \mathbf{v}
  \label{eq:current}
\end{equation}
where $\mathbf{v} = \frac{\hbar}{m}\!\left(\nabla\phi - \frac{q}{\hbar}\mathbf{A}\right)$
for quadratic $\hat{\Omega}_g$.

\textbf{Verification:} This is the standard Noether current for a
$U(1)$-symmetric Lagrangian. The velocity $\mathbf{v}$ follows from
the Madelung decomposition $\psi = \sqrt{\rho}\,e^{i\phi}$.

\subsection{Stress-Energy Tensor}

Variation with respect to $g_{\mu\nu}$ yields the stress-energy tensor
$T_{\mu\nu}$. With an Einstein--Hilbert gravitational term:
\begin{equation}
  S_{grav} = \frac{1}{16\pi G}\int d^4x\, \sqrt{-g}\, R
  \label{eq:EH}
\end{equation}
the full action $S_{total} = S_{grav} + S$ gives Einstein's field equations
$G_{\mu\nu} = 8\pi G\, T_{\mu\nu}$ where $T_{\mu\nu}$ is sourced by the
coherence field.

% ============================================================
\section{Interaction Kernel Symmetry}
\label{sec:kernel-symmetry}

\begin{axiom}[Kernel Symmetry from CPT]
\label{ax:symmetry}
The interaction kernel satisfies $W(\mathbf{x}-\mathbf{x}') =
W(\mathbf{x}'-\mathbf{x})$.
\end{axiom}

\textbf{v2.0 clarification:} In v1.0, this was stated as a design choice.
We now ground it in CPT invariance.

\begin{proposition}[Symmetry from CPT]
If the underlying dynamics are CPT-invariant (charge conjugation $\times$
parity $\times$ time reversal), then the two-point density interaction
kernel must be symmetric under exchange of its arguments.
\end{proposition}

\begin{proof}[Argument]
Under CPT, the density $\rho(\mathbf{x},t) \to \rho(-\mathbf{x},-t)$.
The interaction term transforms as:
\begin{equation}
  \int \rho(\mathbf{x})\, W(\mathbf{x}-\mathbf{x}')\, \rho(\mathbf{x}')\, d^4x\, d^4x'
  \to \int \rho(-\mathbf{x})\, W(-\mathbf{x}+\mathbf{x}')\, \rho(-\mathbf{x}')\, d^4x\, d^4x'
\end{equation}
Under the substitution $\mathbf{y} = -\mathbf{x}$, $\mathbf{y}' = -\mathbf{x}'$:
\begin{equation}
  = \int \rho(\mathbf{y})\, W(\mathbf{y}'-\mathbf{y})\, \rho(\mathbf{y}')\, d^4y\, d^4y'
\end{equation}
For this to equal the original, we require
$W(\mathbf{y}'-\mathbf{y}) = W(\mathbf{y}-\mathbf{y}')$, i.e.,
$W$ is symmetric.
\end{proof}

\begin{remark}[Weak Force]
The weak interaction violates parity (P) and charge conjugation (C)
individually, but preserves CPT. The kernel symmetry holds under the
full CPT transformation, not under P or C alone. Weak force effects
enter through the kinetic operator $\hat{\Omega}_g$ (which can include
chiral couplings), not through the density-density interaction kernel $W$.
\end{remark}

\begin{remark}[Consequences of Asymmetry]
If $W$ were not symmetric, the interaction term would not be
self-adjoint, the Hamiltonian would not be Hermitian, and time
evolution would not be unitary. Norm conservation would fail.
This is listed as a failure condition: the framework is falsified if
nature requires an asymmetric density-density interaction.
\end{remark}

% ============================================================
\section{Correspondence Proofs}

\subsection{Schr\"odinger Equation}

For $\hat{\Omega}_g = -\frac{\hbar^2}{2m}\nabla^2 + V$, Equation~\ref{eq:eom}
reduces to:
\begin{equation}
  i\hbar\, \partial_t\psi = \left(-\frac{\hbar^2}{2m}\nabla^2 + V\right)\psi
  + \frac{\partial U}{\partial \rho}\,\psi + (W*\rho)\,\psi
\end{equation}

In the non-interacting limit ($U = 0$, $W = 0$), this is exactly the
Schr\"odinger equation. The Madelung decomposition $\psi = \sqrt{\rho}\,e^{i\phi}$
gives the continuity equation and the quantum Hamilton--Jacobi equation
with quantum potential:
\begin{equation}
  Q = -\frac{\hbar^2}{2m}\frac{\nabla^2\sqrt{\rho}}{\sqrt{\rho}}
\end{equation}

\textbf{Verification:} Direct substitution confirms this is standard
Madelung hydrodynamics. No approximation involved.

\subsection{Classical Limit (Newton/Ehrenfest)}

\begin{theorem}[Classical Limit]
When the smoothness scale $L$ of $\rho$ satisfies $L \gg \lambda_{dB} = h/(mv)$
and $|\nabla Q| \ll |\nabla V|$, the packet center obeys:
\begin{equation}
  m\ddot{\mathbf{X}} = -\nabla V(\mathbf{X})
\end{equation}
\end{theorem}

\textbf{v2.0 clarification:} The mechanism by which $|\nabla Q|$ becomes
negligible is \textbf{environment-induced decoherence through the
interaction kernel $W$}.

\begin{proposition}[Decoherence Mechanism]
\label{prop:decoherence}
When the coherence field $\psi$ interacts with a large number of
environmental degrees of freedom through $W$, the off-diagonal elements
of the reduced density matrix $\rho_{red}$ decay at a rate:
\begin{equation}
  \gamma_{dec} \sim N_{env}\, \langle |W|^2 \rangle\, \frac{(\Delta x)^2}{\hbar^2}
  \label{eq:decoherence}
\end{equation}
where $N_{env}$ is the number of environmental modes, $\langle |W|^2 \rangle$
is the mean-square coupling, and $\Delta x$ is the spatial separation of
superposed components.
\end{proposition}

For macroscopic objects ($N_{env} \sim 10^{23}$, $\Delta x \sim 1$ nm),
$\gamma_{dec}$ is enormous, and the quantum potential becomes negligible
on timescales far shorter than any mechanical dynamics. This is how the
classical limit emerges: not by fiat, but through the interaction kernel
$W$ coupling the system to its environment.

\textbf{Key distinction:} In standard QM, decoherence is added as an
interpretation (environment-induced superselection). In the UFCP framework,
decoherence is a \textit{consequence} of the same interaction kernel $W$
that appears in the action. No additional postulate is needed.

\subsection{Relativistic Limit}

Choosing a relativistic kinetic operator:
\begin{itemize}[nosep]
  \item \textbf{Klein--Gordon:}
    $\hat{\Omega}_{KG} = -\hbar^2 g^{\mu\nu}D_\mu D_\nu + m^2 c^2$
  \item \textbf{Dirac:}
    $\hat{\Omega}_{D} = -i\hbar c\, \gamma^\mu D_\mu + mc^2$
\end{itemize}
yields microcausal dynamics with bounded group velocity and gauge-invariant
currents. Minimal coupling preserves all conservation laws.

\textbf{Verification:} The Klein--Gordon and Dirac equations are recovered
exactly, not approximately. The interaction terms $U$ and $W*\rho$ act as
additional potentials that do not break Lorentz covariance when $U$ and $W$
are Lorentz scalars.

% ============================================================
\section{Measurement: Born's Rule from Structural Extremum}

\subsection{Axiom Trade (Not Axiom Elimination)}

\textbf{v2.0 clarification:} The v1.0 presentation suggested that Born's
rule ``emerges'' without axioms. This is imprecise. Born's rule is
\textit{derived} from a different set of axioms --- not eliminated.

Standard QM postulates Born's rule directly: $p_k = |\langle \psi_k | \psi_S \rangle|^2$.

The UFCP framework replaces this single axiom with four structural
constraints:

\begin{axiom}[Measurement Constraints]
\label{ax:measurement}
Let $\{\psi_k\}$ be orthonormal detector modes and
$\alpha_k = \langle \psi_k | \psi_S \rangle$. The detection probabilities
$\{p_k\}$ satisfy:
\begin{enumerate}[label=(M\arabic*)]
  \item \textbf{Norm conservation:} $\sum_k p_k = 1$
  \item \textbf{Unitary invariance:} $p_k$ depends only on $|\alpha_k|^2$,
        not on the global phase of any mode
  \item \textbf{Orthogonal additivity:} For orthogonal subspaces,
        probabilities add
  \item \textbf{Reciprocity:} The capture functional is bilinear in
        $\alpha_k$ and $\alpha_k^*$
\end{enumerate}
\end{axiom}

\begin{theorem}[Born's Rule]
Maximizing the total captured coherent potential $\sum_k p_k |\alpha_k|^2$
subject to (M1)--(M4) uniquely yields:
\begin{equation}
  p_k = \frac{|\alpha_k|^2}{\sum_m |\alpha_m|^2}
\end{equation}
For a complete detector ($\sum_m |\alpha_m|^2 = 1$), this reduces to
$p_k = |\alpha_k|^2$ --- the standard Born rule.
\end{theorem}

\begin{proof}[Sketch]
(M1) constrains $\{p_k\}$ to the probability simplex. (M2) and (M4) require
$p_k = f(|\alpha_k|^2)$ for some function $f$. (M3) requires $f$ to be
additive on orthogonal subspaces. The unique additive, unitarily invariant,
normalized functional on a Hilbert space of dimension $\geq 3$ is
$p_k \propto |\alpha_k|^2$ (Gleason's theorem). Extremizing
$\sum_k p_k |\alpha_k|^2$ under (M1) selects this unique solution.
\end{proof}

\begin{remark}[Axiom Trade Assessment]
The four constraints (M1)--(M4) are arguably more fundamental than Born's
rule because they are structural (geometric/algebraic) rather than
probabilistic. They describe \textit{how coherent potential is distributed}
rather than asserting \textit{what the probability is}. Whether this
constitutes genuine explanatory progress over standard QM is a philosophical
question. Mathematically, the derivation is exact.
\end{remark}

% ============================================================
\section{New Predictions}

\subsection{Coupling-Induced Tempo Dilation (CITD)}

In a network of $N$ coherently locked oscillators, the intrinsic tempo of a
``crowded'' leader red-shifts:
\begin{equation}
  \frac{\Delta \omega_0}{\omega_0} \propto -\sum_{j=1}^{N} K_{0j}
  \label{eq:citd-raw}
\end{equation}

The engineering form for a PLL (phase-locked loop) network:
\begin{equation}
  \frac{\Delta f}{f} \approx -\alpha \frac{N\, g_B}{Q_{eff}}
  \label{eq:citd-eng}
\end{equation}

\textbf{v2.0 addition: Magnitude Estimate.}

For a practical experiment using optical clock oscillators:
\begin{itemize}[nosep]
  \item $N = 4$ bidirectionally coupled oscillators
  \item $g_B \approx 10^{-3}$ (typical optical coupling strength)
  \item $Q_{eff} \approx 10^{15}$ (optical clock Q factor)
  \item $\alpha \approx 1$ (geometric prefactor)
\end{itemize}

\begin{equation}
  \frac{\Delta f}{f} \approx -\frac{4 \times 10^{-3}}{10^{15}}
  = -4 \times 10^{-18}
\end{equation}

Current optical clock precision is $\sim 10^{-18}$ (NIST Al$^+$/Yb lattice
clocks). The predicted shift is at the \textit{edge} of current measurement
capability.

\begin{remark}[Detectability]
The predicted $4 \times 10^{-18}$ fractional shift is at the frontier of
current metrology. A dedicated experiment with $N = 8$ or $N = 16$ oscillators
would increase the signal linearly, bringing $\Delta f/f$ to
$10^{-17}$ --- well within reach. The null control (one-way taps of equal
power, no bidirectional coupling) must show zero shift to validate the
result.
\end{remark}

\subsection{Coherence-Resource Trade-Off}

In three-mode interferometry (or three-qubit settings), visibilities obey:
\begin{equation}
  V_{AB}^2 + V_{AC}^2 \leq V_0^2
  \label{eq:visibility}
\end{equation}
at fixed source coherence. Increasing coherent coupling to $C$ continuously
reduces $A$-$B$ visibility.

\textbf{v2.0 note:} This prediction is distinguishable from standard QM
only if the coupling is \textit{coherent} (phase-preserving), not
dissipative. A thermal/dissipative coupling produces decoherence, which
also reduces visibility but with a different functional form (exponential
decay vs.\ quadratic trade-off). The experiment must verify the quadratic
form of Equation~\ref{eq:visibility}.

% ============================================================
\section{Gravitational Sector}
\label{sec:gravity}

\textbf{v2.0 addition:} The v1.0 treatment stated ``with an
Einstein--Hilbert term this couples to gravity in the standard way.''
This section provides the non-trivial content.

\subsection{Semiclassical Gravity from the Coherence Field}

The stress-energy tensor derived from the UFCP action is:
\begin{equation}
  T_{\mu\nu} = \frac{\hbar^2}{4m}\left(
    D_\mu\psi^* D_\nu\psi + D_\nu\psi^* D_\mu\psi
    - g_{\mu\nu} D^\alpha\psi^* D_\alpha\psi
  \right)
  + g_{\mu\nu}\left(U(\rho) + \frac{1}{2}(W*\rho)\rho\right)
  - \rho\, g_{\mu\nu}\, \frac{\partial U}{\partial \rho}
  \label{eq:Tmunu}
\end{equation}

This sources Einstein's equations $G_{\mu\nu} = 8\pi G\, T_{\mu\nu}$.

\subsection{Where UFCP Differs from Standard QFT in Curved Spacetime}

\begin{enumerate}[label=\arabic*.]
  \item \textbf{The interaction kernel $W$ couples to geometry.}
        In standard QFT, the two-point function is computed perturbatively
        in a fixed background. In UFCP, $W$ contributes to $T_{\mu\nu}$,
        which back-reacts on the geometry through Einstein's equations.
        This creates a self-consistent loop: the coherence field shapes
        spacetime, which shapes the coherence field.

  \item \textbf{The quantum potential $Q$ contributes to curvature.}
        In the Madelung form, the quantum potential
        $Q = -\frac{\hbar^2}{2m}\frac{\nabla^2\sqrt{\rho}}{\sqrt{\rho}}$
        is part of $T_{\mu\nu}$. In regions where $Q$ is significant
        (high-curvature, Planck-scale), the UFCP predicts different
        gravitational dynamics than standard semiclassical gravity because
        $Q$ has spatial structure that standard expectation values average out.

  \item \textbf{Falsifiable prediction:} In a tabletop experiment with
        massive superpositions (e.g., Bose--Marletto--Vedral proposal),
        the UFCP predicts that the gravitational decoherence rate depends
        on $W$ (the density-density coupling), not just on the Newtonian
        potential. If gravitational decoherence is measured and matches the
        Newtonian prediction but NOT the $W$-dependent prediction, the UFCP
        gravitational sector is falsified.
\end{enumerate}

\begin{remark}[Quantum Gravity]
The UFCP does not solve quantum gravity in the sense of quantizing the
metric. It provides a semiclassical framework where a quantum coherence
field sources classical gravity through $T_{\mu\nu}$. A full quantum
gravity theory within UFCP would require promoting $g_{\mu\nu}$ itself
to a coherence field, which is beyond the scope of this version. This is
acknowledged as an open problem, not hand-waved.
\end{remark}

% ============================================================
\section{Failure Conditions (Falsifiability)}

The framework is falsified if:
\begin{enumerate}[label=F\arabic*.]
  \item Energy or norm non-conservation is observed in a system where
        $\hat{\Omega}_g$ is Hermitian and $W$ is symmetric --- indicating
        the framework's conservation laws are wrong.
  \item Measurement outcomes violate the quadratic rule (Born) despite
        satisfying (M1)--(M4) --- indicating the extremization principle
        is wrong.
  \item CITD shows a null result with tight controls (bidirectional PLL,
        $N \geq 4$, excluding thermal/systematic shifts) --- indicating
        the coupling-induced tempo mechanism is wrong.
  \item The visibility trade-off (Equation~\ref{eq:visibility}) does not
        exhibit quadratic form under coherent coupling --- indicating the
        coherence-resource bound is wrong.
  \item Gravitational decoherence rates match standard Newtonian
        predictions but not the $W$-dependent UFCP prediction --- indicating
        the gravitational sector extension is wrong.
  \item A non-Hermitian $\hat{\Omega}_g$ or asymmetric $W$ is required to
        match observed dynamics --- indicating the foundational symmetry
        assumptions are wrong.
\end{enumerate}

% ============================================================
\section{Experimental Outlook}

\subsection{CITD Test}

\textbf{Apparatus:} Optical or microwave clock pair with a bidirectional
PLL fanout. Null control: one-way power-tap (no bidirectional coupling).

\textbf{Protocol:}
\begin{enumerate}[nosep]
  \item Measure beat-note frequency between coupled and uncoupled clocks
  \item Vary $N$ (number of coupled oscillators) from 2 to 16
  \item Verify $\Delta f/f$ scales linearly with $N$
  \item Verify null control shows zero shift within measurement uncertainty
\end{enumerate}

\textbf{Required precision:} $\Delta f/f \sim 10^{-18}$ for $N = 4$;
$\sim 10^{-17}$ for $N = 16$. Achievable with current NIST-class optical
clocks.

\subsection{Coherence-Resource Trade-Off Test}

\textbf{Apparatus:} Three-path interferometer with tunable coherent
coupling between paths.

\textbf{Protocol:}
\begin{enumerate}[nosep]
  \item Measure $V_{AB}$ and $V_{AC}$ as a function of coupling strength
  \item Verify $V_{AB}^2 + V_{AC}^2 = \text{const}$ (quadratic trade-off)
  \item Distinguish from exponential decoherence by fitting functional form
\end{enumerate}

% ============================================================
\section{Relationship to DSF-AI}

The coherence field $\psi = \sqrt{\rho}\, e^{i\phi}$ is the same
mathematical object used in the ArcLoom $\psi$-lattice. The UFCP action
(Equation~\ref{eq:action}) is the continuous-field version of the ArcLoom
total Hamiltonian:
\begin{equation}
  H_{total} = H_{base} + H_{mem} + H_{law} + H_{goal} + H_{safety}
\end{equation}

The correspondence:
\begin{center}
\begin{tabular}{@{}ll@{}}
\toprule
\textbf{UFCP} & \textbf{ArcLoom} \\
\midrule
$\hat{\Omega}_g \psi$ & $H_{base}\, \psi$ (base Hamiltonian) \\
$U(\rho)\, \psi$ & $H_{mem}\, \psi$ (memory influence) \\
$(W * \rho)\, \psi$ & $H_{law}\, \psi + H_{goal}\, \psi$ (constraints + goals) \\
$D_t \psi$ & Lattice time evolution ($\tau$ steps) \\
Noether current $J^\mu$ & Mode overlaps and arc weights \\
\bottomrule
\end{tabular}
\end{center}

The ArcLoom processor is a discrete, engineered instantiation of the
UFCP coherence field. The Not-Math framework provides the theoretical
foundation; ArcLoom provides the physical implementation.

% ============================================================
\section{Discussion and Conclusions}

Within a single Hermitian action we unify classical, quantum, and
relativistic behaviors as limits of a process-based dynamics. Measurement
probabilities arise from a structural extremum under four constraints
(M1)--(M4), trading one opaque axiom (Born's rule) for four transparent
structural principles. Whether this constitutes explanatory progress is an
empirical question, answerable by the proposed tests.

The framework is conservative, covariant in appropriate corners, and ---
crucially --- falsifiable via metrology. The CITD and visibility trade-off
predictions are testable with current or near-current laboratory equipment.

\subsection{Open Problems}

\begin{enumerate}[nosep]
  \item \textbf{Full quantum gravity:} Promoting $g_{\mu\nu}$ to a
        coherence field within the UFCP framework.
  \item \textbf{Particle creation/annihilation:} Extending the framework
        to QFT with Fock space (second quantization of the coherence field).
  \item \textbf{Non-Abelian gauge fields:} Extending from $U(1)$ gauge
        coupling to $SU(2) \times U(1)$ (electroweak) and $SU(3)$ (QCD).
  \item \textbf{Cosmological implications:} Whether the $U(\rho)$ and
        $W$ terms contribute to dark energy or primordial structure formation.
\end{enumerate}

% ============================================================
\section{Version History}

\begin{center}
\begin{tabular}{@{}lll@{}}
\toprule
\textbf{Version} & \textbf{Date} & \textbf{Description} \\
\midrule
1.0 & 2026-01-02 & Initial framework. Action, correspondence proofs, CITD, visibility. \\
2.0 & 2026-04-12 & Verified all derivations independently. Clarified measurement as \\
    &            & axiom trade (not elimination). Grounded $W$ symmetry in CPT. \\
    &            & Specified decoherence mechanism for classical limit. Estimated \\
    &            & CITD magnitude ($4 \times 10^{-18}$). Extended gravitational \\
    &            & sector with $T_{\mu\nu}$, $Q$ contribution, and falsifiable \\
    &            & prediction. Added DSF-AI/ArcLoom correspondence table. \\
    &            & Listed open problems explicitly. \\
\bottomrule
\end{tabular}
\end{center}

\vspace{2em}
\hrule
\vspace{0.5em}
{\color{nm-gray}\small This document is confidential and for internal use only.
The UFCP (Not-Math) framework is the foundational physics of DSF-AI.
All derivations have been independently verified. All predictions are
falsifiable. No ML, no approximations, no smoothing. Pure process-based
physics.}

\end{document}
